\section{Introducción}

El presente trabajo práctico detalla el diseño, simulación e implementación física de un sistema de muestreo y acondicionamiento de señales. Se aborda el desarrollo de las etapas analógicas, comenzando por un oscilador de relajación empleado para el control del muestreo y la sincronización de conmutadores. Posteriormente, se analizan e implementan filtros activos pasabajos basados en topologías Sallen-Key con aproximación de Chebyshev de quinto orden, destinados a operar como filtro anti-alias (FAA) y filtro recuperador (FR).

Se evalúa el funcionamiento de la etapa de Sample \& Hold mediante el circuito integrado LF398, estudiando las no idealidades del componente, los tiempos de adquisición y la selección óptima del capacitor de retención. Además, se documenta el diseño de la placa de circuito impreso (PCB) aplicando criterios de ruteo orientados a la reducción de ruido, separación de dominios analógico-digitales y protección de nodos de alta impedancia. En la instancia final, se realiza un análisis teórico y empírico comparando el muestreo instantáneo, natural y sub-Nyquist bajo distintas condiciones experimentales.